\section{Ideal models on the actual variety product}

\begin{definition}[The actual polynomial-line product of the variety]
\label{def:pl-native-variety-product}
\uses{prop:pl-vertical-ideal-models,projective:actual-product-model-ideal-functions}
\lean{EnvelopeProperty.ProjSubvariety.nativeHomogeneousIdeal,EnvelopeProperty.ProjSubvariety.nativeIdeal,EnvelopeProperty.ProjSubvariety.NativeScheme,EnvelopeProperty.ProjSubvariety.nativeProductCutIdeal,EnvelopeProperty.ProjSubvariety.NativeProductModel,EnvelopeProperty.ProjSubvariety.nativeProductModelIso,EnvelopeProperty.ProjSubvariety.nativeProductModelIso_over_base}
\leanok
Let $X\subset\mathbf P^M_k$ be the given homogeneous projective subvariety.
Its native scheme is the actual closed subscheme cut out by its original
homogeneous ideal.  In $S=\mathbf P^M_k\times\mathbf A^1_k$, pull that
ideal sheaf back by the first projection, and denote the resulting closed
subscheme by $S_X$.  There is an actual whole-scheme isomorphism
\[
 S_X\simeq X\times_{\operatorname{Spec}k}\operatorname{Spec}k[t].
\]
\end{definition}
\begin{proof}
The closed subscheme of a pulled-back ideal is isomorphic to the fibre
product of its original closed immersion with the pullback morphism.
Apply this to the first projection $S\to\mathbf P^M_k$.  Symmetry of the
fibre product and the associativity isomorphism for two consecutive fibre
products identify this scheme with
$X\times_{\operatorname{Spec}k}\operatorname{Spec}k[t]$.
The isomorphism preserves the actual polynomial-base structure morphism.
\end{proof}

\begin{lemma}[Actual Gauss points of the variety product]
\label{lem:pl-native-variety-gauss-lift}
\uses{def:pl-native-variety-product}
\lean{EnvelopeProperty.ProjSubvariety.gaussProductPoint_mem_nativeProductCutIdeal,EnvelopeProperty.ProjSubvariety.nativeProductGaussPoint,EnvelopeProperty.ProjSubvariety.map_nativeProductGaussPoint,EnvelopeProperty.ProjSubvariety.continuous_nativeProductGaussPoint}
\leanok
Every $v\in X^{\rm an}$ has an actual valued scheme point of $S_X$ above
the prescribed polynomial Gauss seminorm, whose entire image in $S$ is
the ambient Gauss point $\sigma(v)$.  This point varies continuously with
$v$.
\end{lemma}
\begin{proof}
Coefficient base change identifies the pulled-back defining ideal with
the homogeneous ideal generated by the images of the original equations
of $X$ in $k[t][x_0,\ldots,x_M]$.  The normalized Gauss seminorm kills
each such image: coefficient interchange makes it the constant
$t$-polynomial associated to an original equation, and its value is zero
because $v\in X^{\rm an}$.  Its kernel is an ideal, so it kills the
entire generated ideal.  The actual support therefore belongs to the
closed subscheme's support.  The closed-immersion valued-point
homeomorphism supplies a unique lift preserving the full residue-field
valuation and the polynomial-base point.  Its inverse is continuous,
which proves the continuity assertion.
\end{proof}

\begin{lemma}[Whole intrinsic ideal logarithms survive restriction]
\label{lem:pl-native-variety-ideal-log}
\uses{lem:pl-native-variety-gauss-lift,prop:pl-vertical-ideal-models}
\lean{EnvelopeProperty.ProjSubvariety.nativeProductIdeal,EnvelopeProperty.ProjSubvariety.modelIdealLog_nativeProductIdeal,EnvelopeProperty.ProjSubvariety.IsPL.exists_nativeProductIdeal_models,EnvelopeProperty.ProjSubvariety.IsLocallyConvex.exists_nativeProductIdeal_models}
\leanok
For every actual homogeneous ambient product ideal $J$, restrict its
actual ideal sheaf to $S_X$.  Its entire intrinsic model logarithm at the
constructed Gauss lift equals its original ambient model logarithm.
Consequently the exact finite logarithm-difference formula for every PL
function holds for genuine ideals on $X\times\mathbf A^1$.
\end{lemma}
\begin{proof}
The original defining inclusion is a closed immersion, hence proper;
composing with the proper ambient product structure gives a proper
structure morphism on $S_X$.  Functoriality of intrinsic ideal norms
under actual geometric ideal pullback identifies the norm of the
restricted ideal at a source point with the norm of the original ideal
at its entire image.  Apply this identity to the constructed Gauss
point, whose image is the original ambient point, and take logarithms.
Apply the equality separately to the two constructed PL ideals.
\end{proof}
